%% kapa: exam=1 round=24 year=2013 period=2 form=A subject=accounting subject_ko=회계학 seq=18 number=18
%% source: Q-Net 감정평가사 기출문제 (cst00309) artlSeq=5073802 file=extracted/1008/5073802/2013년 제24회 감정평가사 1차시험 기출문제/2교시(A형).hwp
%% conversion: kordoc 4.24.51 -> tools/gen1.py (auto); review=proofread-sonnet; flags=none
\begin{question}{18}
(주)대한은 다음의 사채를 사채권면에 표시된 발행일(20×1년 1월 1일)이 아닌 20×1년 4월 1일에 실제 발행하였다.
\begin{framedtext}
○ 만 기 일: 20×3년 12월 31일 \par
○ 액면금액: ￦100,000 \par
○ 표시이자율: 연 10\% \par
○ 이자는 매년 말에 지급한다.
\end{framedtext}
20×1년 4월 1일 (주)대한의 시장이자율이 연 12\%일 경우, 20×1년 4월 1일의 사채발행이 동 시점의 (주)대한의 부채총액에 미치는 영향은? (단, 현가계수는 아래의 표를 이용하며, 이자는 월할계산한다. 단수차이로 인한 오차가 있으면 가장 근사치를 선택한다.) \par
\uline{현가계수표} \par
{\small\centering
\begin{tabular}{|c|c|c|}
\hline
\makecell{(3년 기준)} & \makecell{10\%} & \makecell{12\%} \\
\hline
\makecell{단일금액 ￦1의 현가계수} & \makecell{0.7513} & \makecell{0.7118} \\
\hline
\makecell{정상연금 ￦1의 현가계수} & \makecell{2.4868} & \makecell{2.4018} \\
\hline
\end{tabular}\par}

\begin{choices}
  \item ￦95,554 증가
  \item ￦97,698 증가
  \item ￦98,054 증가
  \item ￦100,000 증가
  \item ￦102,500 증가
\end{choices}
\end{question}
\answer{3}
